\subsubsection{Mixed products of rank-one forms on three rows}

\begin{lemma}\label{B:three-row-rank-one-products}
Suppose the minimal ideal of a finite monoid is a rectangular
band with three rows.  If two rational degree-two classes have
rank-one balanced matrices, their product is zero.  The same
statement holds with rows and columns interchanged.  No finiteness
assumption on the homotopy type is required.
\end{lemma}
\begin{proof}
Write the two balanced forms as $L_i=b_i\otimes a_i$, for $i=1,2$.
Use the affine representatives and notation of
Lemma~\ref{B:rank-one-square}: normalize the row-value functions
$a_i(r)$ to vanish at a distinguished row, and similarly normalize
the column-value functions.  There are rational characters $\chi_i$
such that
\[
 a_i(\sigma_m r)=a_i(m)+\chi_i(m)a_i(r),\qquad
 b_i(\tau_m c)=b_i(m)+\chi_i(m)b_i(c),
\]
and the class of $L_i$ is represented by
$u_i(m,n)=b_i(m)a_i(n)$.  The corresponding multiplication identities
are
\[
 a_i(mn)=a_i(m)+\chi_i(m)a_i(n),\qquad
 b_i(mn)=b_i(n)+\chi_i(n)b_i(m).
\]

If the two row functionals are dependent, every linear combination
of $L_1,L_2$ has rank at most one.  Their individual squares and
the square of their sum vanish by Lemma~\ref{B:rank-one-square};
polarization over $\mathbb Q$ gives $[u_1][u_2]=0$.

Otherwise the three functions $1,a_1,a_2$ form a basis of the
functions on the three-point row set.  Since both $a_i$ vanish at
the distinguished row, there are rational numbers $s,t$ with
\[
 a_1(r)a_2(r)=s a_1(r)+t a_2(r)
 \quad\text{for every row }r.
\]
Apply this identity to $\sigma_m r$ and substitute the affine
formulas.  The coefficient of $a_2(r)$ is
\[
 \chi_2(m)\bigl(a_1(m)+t(\chi_1(m)-1)\bigr)=0.
\]
Set $a'_1(m)=a_1(m)+t(\chi_1(m)-1)$.  This is still an affine
cocycle with character $\chi_1$, and
\[
 u'_1(m,n):=b_1(m)a'_1(n),\qquad
 u'_1-u_1=-t\delta b_1.
\]
Define a rational three-cochain by
\[
 w(m,n,p)=b_1(m)a'_1(n)b_2(n)a_2(p).
\]
Expanding the bar coboundary using the two multiplication
identities gives the exact equality
\begin{align*}
 (\delta w)(m,n,p,q)
 &=u'_1(m,n)u_2(p,q)\\
 &\quad+b_1(m)b_2(n)\chi_1(n)
              \chi_2(p)a'_1(p)a_2(q)\\
 &=u'_1(m,n)u_2(p,q).
\end{align*}
The last equality uses $\chi_2(p)a'_1(p)=0$.  Thus the mixed
product is a coboundary.  Applying the result to the opposite
monoid proves the column version.
\end{proof}

\begin{corollary}\label{B:three-row-semisimple-cups}
Let $M$ have a three-row rectangular-band minimal ideal, and let
$V$ be its two-dimensional rational row augmentation module.
If $V$ is a direct sum of one-dimensional modules, every product
of two rational degree-two classes of $BM$ is zero.
\end{corollary}
\begin{proof}
An equivariant decomposition of $V$ into its two one-dimensional
summands decomposes every balanced map $V\to W^*$ into two
balanced maps of rank at most one.  All pairwise products of
these summands vanish by Lemma~\ref{B:three-row-rank-one-products}.
\end{proof}

\begin{lemma}\label{B:six-dimensional-simple-row}
Let $A$ be a rational algebra of dimension at most six, and let
$V$ be a simple left $A$-module of rational dimension two.
For every right module $W$, either $V$ is projective or
$\dim_{\mathbb Q}\operatorname{Tor}^A_n(W,V)$ is even for every
$n$.  In particular, a one-dimensional degree-two Tor group
cannot occur.
\end{lemma}
\begin{proof}
Write $D=\operatorname{End}_A(V)$.  The division algebra $D$ has
rational dimension one or two.  In the second case every Tor
group is a module over $D$ and has even rational dimension.
In the first case the corresponding simple factor of
$A/J$, where $J=\operatorname{rad}A$, is $M_2(\mathbb Q)$.
Every nonzero bimodule joining this factor to another simple
factor has rational dimension at least two; a bimodule from
this factor to itself has dimension at least four.

If this is the only factor of $A/J$, then $\dim J\leq2$ makes
its component in $J/J^2$ zero, and nilpotence gives $J=0$.
Otherwise $\dim(A/J)\geq5$ and $\dim J\leq1$, so again no
component of $J/J^2$ can meet the $M_2(\mathbb Q)$ factor.
Products of radical components show that no component of any
$J^n/J^{n+1}$ meets it either.  Hence this matrix factor splits
off as a semisimple direct factor of $A$, and its simple module
$V$ is projective.
\end{proof}

\begin{corollary}\label{B:three-by-six-order24-row-reduction}
For a possible twenty-four-element monoid model of
$S^2\times S^2$ whose minimal ideal is a $3\times6$ band,
the two-dimensional row augmentation module over the
six-dimensional contracted rational algebra must be
indecomposable and nonsemisimple.  Its image algebra in
$\operatorname{End}_{\mathbb Q}(V)$ is either
$\mathbb Q[\epsilon]/(\epsilon^2)$ or the algebra of upper
triangular two-by-two rational matrices.
\end{corollary}
\begin{proof}
The band Tor formula gives $\dim\operatorname{Tor}_2(W,V)=1$.
The preceding lemma excludes a two-dimensional simple $V$,
and Corollary~\ref{B:three-row-semisimple-cups} excludes a sum
of one-dimensional simples.  A decomposable two-dimensional
module would have precisely that latter form.  Thus $V$ is a
nonsplit extension of two rational characters.  Choose a basis
adapted to its one-dimensional submodule.  If the characters
are equal, the image algebra is the dual-number algebra; if
they are distinct, it contains the two diagonal idempotents
and a nonzero off-diagonal matrix, and is the full triangular
algebra.
\end{proof}

\begin{proposition}\label{B:three-row-dual-cup-radical}
Suppose the minimal ideal of a finite monoid $M$ is a rectangular
band with three rows, and its rational row augmentation module
$V$ has image algebra $\mathbb Q[\epsilon]/(\epsilon^2)$ acting
nonsemisimply.  Then either all products of rational degree-two
classes vanish, or their cup-product pairing has a nonzero common
radical.  Consequently $BM$ cannot have the cohomology ring of
$S^2\times S^2$.  If $\dim H^2(BM;\mathbb Q)=1$, its generator
has square zero as well.
\end{proposition}
\begin{proof}
First identify the actual row maps, including their affine data.
The image algebra contains a nonzero nilpotent matrix.  Some row
actor must have a nonscalar augmentation operator, which in a
dual-number algebra has the form $\lambda I+\mu N$ with
$\mu\ne0$ and $N^2=0$.  Its powers belong to a finite set.
In characteristic zero this forces $\lambda=0$: a nonzero
$\lambda$ would make the nontrivial Jordan part have infinitely
many powers.  Thus this actor is a three-point chain.  After
relabeling its map is $(0,0,1)$, and in the augmentation basis
$(e_1-e_0,e_2-e_0)$ its matrix is
\[
 N=\begin{pmatrix}0&1\\0&0\end{pmatrix}.
\]
A three-point map whose augmentation operator belongs to
$\operatorname{span}(I,N)$ is one of the following six maps:
\[
 (0,1,2),\quad (0,0,1),\quad(1,1,0),\quad
 (0,0,0),\quad(1,1,1),\quad(2,2,2).
\]
This list follows directly by equating the two columns
$e_{\sigma(1)}-e_{\sigma(0)}$ and
$e_{\sigma(2)}-e_{\sigma(0)}$ with those of $\lambda I+\mu N$.
Their augmentation operators are respectively $I,N,-N,0,0,0$.

Let $a(m)$ be the two-coordinate vector
$e_{\sigma_m(0)}-e_0$, and let $x$ be the row functional with
$x(e_1-e_0)=0$, $x(e_2-e_0)=1$.  Write
$\alpha(m)=x(a(m))$ and let $\chi(m)$ be the common diagonal
character of the augmentation matrix $A_m$.
Inspection of the six maps gives the exact identity
\[
 \alpha(m)A_m=0.
\]
Indeed $\alpha(m)$ is nonzero only for the constant map to row
$2$.  Also
$\alpha(mn)=\alpha(m)+\chi(m)\alpha(n)$.

Every rank-one balanced form has row factor $x$: the nonzero
nilpotent $N$ has a unique invariant line, so a rank-one
equivariant map from $V$ has kernel $\operatorname{im}N$.
Write its affine cocycle as
$v(m,n)=r(m)\alpha(n)$, where
$r(mn)=r(n)+\chi(n)r(m)$.
An arbitrary balanced form has affine cocycle
$u(m,n)=b(m)a(n)$, where $b(m)$ is a row vector and
\[
 a(mn)=a(m)+A_m a(n),\qquad
 b(mn)=b(n)+b(m)A_n.
\]
Define
\[
 w(m,n,p)=r(m)\alpha(n)u(n,p).
\]
Using the cocycle identity for $u$ and the displayed affine
identities, the full coboundary is
\begin{align*}
 (\delta w)(m,n,p,q)
 &=v(m,n)u(p,q)\\
 &\quad+r(m)\chi(n)\alpha(p)b(n)A_p a(q)\\
 &=v(m,n)u(p,q).
\end{align*}
Thus every rank-one balanced class annihilates every
degree-two class.

Finally, $N$ commutes with the action algebra.  If a balanced
map $L:V\to W^*$ has rank two, then $LN$ is a nonzero balanced
map of rank one.  The natural identification of balanced forms
with $H^2$ makes this a nonzero cohomology class in the common
radical.  If no balanced map has rank two, every class has rank
at most one and the preceding coboundary formula makes all
products zero.  These are all rank cases.  A one-dimensional
space of balanced forms cannot contain a rank-two map, since
$L$ and $LN$ would be linearly independent; this proves the
last assertion.
\end{proof}

\begin{corollary}\label{B:three-by-six-order24-triangular}
In the situation of
Corollary~\ref{B:three-by-six-order24-row-reduction}, the row
augmentation image is the full triangular two-by-two rational
matrix algebra.  Thus the possible twenty-four-element product
case has a nonsplit extension of two distinct rational characters
as its row augmentation module.
\end{corollary}
\begin{proof}
The preceding proposition excludes the remaining dual-number
alternative, since the degree-two pairing of $S^2\times S^2$
has zero radical.
\end{proof}
